% GENERATED FILE. Edit canonical lessons, then run npm run generate:books.
% canonical-source-hash: fnv1a64:ec5dd952e7ad82ae
% canonical-lessons: BN-C160-buddhiman, BN-C160-porishromi, BN-C160-lobhi, BN-C160-chintito, BN-C160-buro

\chapter{Intelligent, Hardworking, Greedy, Worried, Elderly}
\label{ch:bn-intelligent-hardworking-greedy-worried-elderly}
\hlchaptermodality{160}

\begin{chapteropening}
\textbf{By the end of this chapter:} \emph{I can say intelligent, hardworking, greedy, worried and elderly.}

Complete the last lesson of chapter 160: I can say intelligent, hardworking, greedy, worried and elderly.
\end{chapteropening}

\section[\texorpdfstring{\bn{বুদ্ধিমান} (buddhimān)}{বুদ্ধিমান (buddhimān)}]{\bn{বুদ্ধিমান} (buddhimān) — intelligent}

\label{lesson:BN-C160-buddhiman}



\begin{quote}
Before the new one: say the Bengali for yet, then the Bengali for otherwise.
\end{quote}

\subsection*{You'll want to know: \bn{বুদ্ধিমান}}
\textbf{\bn{বুদ্ধিমান}} — \emph{buddhimān} — \textquotedblleft{}intelligent\textquotedblright{}.

\begin{cousinweb}
\textbf{\bn{বুদ্ধিমান}} (\emph{buddhimān}) is Sanskrit \emph{buddhimat}, \textquotedblleft{}intelligent\textquotedblright{}, taken over whole as a learned loan, a tatsama.
\end{cousinweb}

\begin{grammarlens}[title={What you've built}]
One more word for this chapter.
\end{grammarlens}

\subsection*{Your turn}
\begin{itemize}
\raggedright
  \item \emph{Say it:} \emph{buddhimān}
  \item \emph{Say it:} \emph{buddhimān}, once more
  \item \emph{Say it:} say \emph{buddhimān}
  \item \emph{Recall:} say the Bengali for yet, then the Bengali for otherwise, then say \emph{buddhimān} again
\end{itemize}

\begin{tcolorbox}[breakable,colback=teal!4,colframe=teal!35!black,title={Before you move on}]
What does \bn{বুদ্ধিমান} mean? (\textquotedblleft{}Intelligent\textquotedblright{}.) Say it once more.
\end{tcolorbox}
\section[\texorpdfstring{\bn{পরিশ্রমী} (pôrishrômī)}{পরিশ্রমী (pôrishrômī)}]{\bn{পরিশ্রমী} (pôrishrômī) — hardworking}

\label{lesson:BN-C160-porishromi}



\begin{quote}
Before the new one: say the Bengali for otherwise, then the Bengali for intelligent.
\end{quote}

\subsection*{You'll want to know: \bn{পরিশ্রমী}}
\textbf{\bn{পরিশ্রমী}} — \emph{pôrishrômī} — \textquotedblleft{}hardworking\textquotedblright{}.

\begin{cousinweb}
\textbf{\bn{পরিশ্রমী}} (\emph{pôrishrômī}) is Sanskrit \emph{pariśrama}, \textquotedblleft{}exertion, toil\textquotedblright{}, taken over whole, + the ending \emph{-ī}, \textquotedblleft{}having\textquotedblright{}: \textquotedblleft{}having toil\textquotedblright{}.
\end{cousinweb}

\begin{grammarlens}[title={What you've built}]
One more word for this chapter.
\end{grammarlens}

\subsection*{Your turn}
\begin{itemize}
\raggedright
  \item \emph{Say it:} \emph{pôrishrômī}
  \item \emph{Say it:} \emph{pôrishrômī}, once more
  \item \emph{Say it:} say \emph{pôrishrômī}
  \item \emph{Recall:} say the Bengali for otherwise, then the Bengali for intelligent, then say \emph{pôrishrômī} again
\end{itemize}

\begin{tcolorbox}[breakable,colback=teal!4,colframe=teal!35!black,title={Before you move on}]
What does \bn{পরিশ্রমী} mean? (\textquotedblleft{}Hardworking\textquotedblright{}.) Say it once more.
\end{tcolorbox}
\section[\texorpdfstring{\bn{লোভী} (lobhī)}{লোভী (lobhī)}]{\bn{লোভী} (lobhī) — greedy}

\label{lesson:BN-C160-lobhi}



\begin{quote}
Before the new one: say the Bengali for intelligent, then the Bengali for hardworking.
\end{quote}

\subsection*{You'll want to know: \bn{লোভী}}
\textbf{\bn{লোভী}} — \emph{lobhī} — \textquotedblleft{}greedy\textquotedblright{}.

\begin{cousinweb}
\textbf{\bn{লোভী}} (\emph{lobhī}) is Sanskrit \emph{lobhin}, \textquotedblleft{}greedy\textquotedblright{}, taken over whole as a learned loan, a tatsama (CDIAL 11147). It is built on \emph{lobha}, \textquotedblleft{}eager desire\textquotedblright{}, which CDIAL lists.
\end{cousinweb}

\begin{grammarlens}[title={What you've built}]
One more word for this chapter.
\end{grammarlens}

\subsection*{Your turn}
\begin{itemize}
\raggedright
  \item \emph{Say it:} \emph{lobhī}
  \item \emph{Say it:} \emph{lobhī}, once more
  \item \emph{Say it:} say \emph{lobhī}
  \item \emph{Recall:} say the Bengali for intelligent, then the Bengali for hardworking, then say \emph{lobhī} again
\end{itemize}

\begin{tcolorbox}[breakable,colback=teal!4,colframe=teal!35!black,title={Before you move on}]
What does \bn{লোভী} mean? (\textquotedblleft{}Greedy\textquotedblright{}.) Say it once more.
\end{tcolorbox}
\section[\texorpdfstring{\bn{চিন্তিত} (chintitô)}{চিন্তিত (chintitô)}]{\bn{চিন্তিত} (chintitô) — worried}

\label{lesson:BN-C160-chintito}



\begin{quote}
Before the new one: say the Bengali for hardworking, then the Bengali for greedy.
\end{quote}

\subsection*{You'll want to know: \bn{চিন্তিত}}
\textbf{\bn{চিন্তিত}} — \emph{chintitô} — \textquotedblleft{}worried\textquotedblright{}.

\begin{cousinweb}
\textbf{\bn{চিন্তিত}} (\emph{chintitô}) is Sanskrit \emph{cintita}, \textquotedblleft{}thought about\textquotedblright{}, taken over whole as a learned loan, a tatsama.
\end{cousinweb}

\begin{grammarlens}[title={What you've built}]
One more word for this chapter.
\end{grammarlens}

\subsection*{Your turn}
\begin{itemize}
\raggedright
  \item \emph{Say it:} \emph{chintitô}
  \item \emph{Say it:} \emph{chintitô}, once more
  \item \emph{Say it:} say \emph{chintitô}
  \item \emph{Recall:} say the Bengali for hardworking, then the Bengali for greedy, then say \emph{chintitô} again
\end{itemize}

\begin{tcolorbox}[breakable,colback=teal!4,colframe=teal!35!black,title={Before you move on}]
What does \bn{চিন্তিত} mean? (\textquotedblleft{}Worried\textquotedblright{}.) Say it once more.
\end{tcolorbox}
\section[\texorpdfstring{\bn{বুড়ো} (buṛo)}{বুড়ো (buṛo)}]{\bn{বুড়ো} (buṛo) — elderly, old (of a person)}

\label{lesson:BN-C160-buro}



\begin{quote}
Before the new one: say the Bengali for greedy, then the Bengali for worried.
\end{quote}

\subsection*{You'll want to know: \bn{বুড়ো}}
\textbf{\bn{বুড়ো}} — \emph{buṛo} — \textquotedblleft{}elderly, old (of a person)\textquotedblright{}.

\begin{cousinweb}
\textbf{\bn{বুড়ো}} (\emph{buṛo}) is a tadbhava, a word inherited and reshaped by sound change. It comes from \emph{*buḍḍha}, \textquotedblleft{}old\textquotedblright{}, a form CDIAL reconstructs from its descendants (CDIAL 9271). CDIAL lists Bengali \emph{buṛā}, \textquotedblleft{}old\textquotedblright{}, and notes that the initial \emph{b-} makes direct descent from Sanskrit \emph{vṛddha}, \textquotedblleft{}grown\textquotedblright{}, unlikely.
\end{cousinweb}

\begin{grammarlens}[title={What you've built}]
One more word for this chapter.
\end{grammarlens}

\subsection*{Your turn}
\begin{itemize}
\raggedright
  \item \emph{Say it:} \emph{buṛo}
  \item \emph{Say it:} \emph{buṛo}, once more
  \item \emph{Say it:} say \emph{buṛo}
  \item \emph{Recall:} say the Bengali for greedy, then the Bengali for worried, then say \emph{buṛo} again
\end{itemize}

\begin{tcolorbox}[breakable,colback=teal!4,colframe=teal!35!black,title={Before you move on}]
What does \bn{বুড়ো} mean? (\textquotedblleft{}Elderly, old (of a person)\textquotedblright{}.) Say it once more.
\end{tcolorbox}
